\begin{solapartat}
\figura[0.4]{sol_fig1.png}

\punts{0,35} Per simetria serà el punt central del triangle (baricentre). És equidistant
a les tres càrregues, per tant, el mòdul dels camps creats serà igual.
\[ E_1 = E_2 = E_3 = k\,\frac{q}{r^2} \]

La component $x$ dels camps $\vec{E}_2$ i $\vec{E}_3$ es cancel·len mútuament.

En aquest cas, l'angle que formen els camps respecte la vertical és $60^\circ$, i la
component $y$ de la suma dels camps $\vec{E}_2$ i $\vec{E}_3$ és:
% DUBTE: a l'original el denominador és "4^2" en comptes de "r^2"; sembla una errata,
% es transcriu tal com surt a la pauta oficial.
\[ E_{23,y} = -E_2\cos(60°) - E_3\cos(60°) = -2k\,\frac{q}{4^2}\,\frac{1}{2} = -k\,\frac{q}{r^2} \]

I, finalment, el camp total és:
\[ E_{T,y} = E_1 - E_{23,y} = k\,\frac{q}{r^2} - k\,\frac{q}{r^2} = 0 \]

La justificació també es pot basar en arguments exclusivament geomètrics i de simetria.
Per exemple, atès que les tres càrregues són iguals i el triangle és equilàter, si
suposem que el camp elèctric en el centre no és nul, només cal fer una rotació de
$120^\circ$ respecte el punt central (una rotació de $120^\circ$ dona una distribució de
càrrega idèntica) i veurem que el mateix problema té dues solucions diferents, cosa que
és impossible, per tant, el camp en aquest punt ha de ser nul.

\figura[0.32]{sol_fig2.png}

\punts{0,4} Calculem el camp al vèrtex superior. Per simetria sabem que el resultat no
depèn del vèrtex triat.

La direcció del camp elèctric creat per una càrrega correspon a la línia que uneix la
càrrega amb el vèrtex on calculem el camp, és a dir, el camp està dirigit segons el
costat del triangle que uneix els dos vèrtexs. Com que les dues càrregues de la base són
negatives, el sentit apuntarà cap a la càrrega que crea el camp, com s'indica al dibuix.

\punts{0,1} El camp total serà la suma vectorial del camp creat per les dues càrregues:
\[ \vec{E}_T = \vec{E}_1 + \vec{E}_2 \]

\punts{0,2} Com que les dues càrregues són iguals i estan a la mateixa distància, el
mòdul del camp que creen és el mateix. $E_1 = E_2 = k\,\dfrac{q}{d^2}$

\punts{0,2} Els camps $\vec{E}_1$ i $\vec{E}_2$ formen el mateix angle respecte la
vertical, que és la meitat de l'angle que formen els dos costats del triangle. Com que és
un triangle equilàter, els angles del triangle són de $60^\circ$ i l'angle que formen els
camps respecte la vertical serà $30^\circ$.

Llavors les components $x$ es cancel·len perquè són d'igual magnitud però tenen sentits
oposats, mentre que les components $y$ es sumen perquè tenen el mateix sentit (negatiu
perquè apunten cap avall):
\[ E_{T,x} = E_2\sin(30°) - E_1\sin(30°) = 0 \]
\[ E_{T,y} = -E_2\cos(30°) - E_2\cos(30°) = -2k\,\frac{q}{d^2}\,\frac{\sqrt{3}}{2} = -\sqrt{3}\,k\,\frac{q}{d^2} \]

Així doncs, el camp és: $\vec{E}_T = -\sqrt{3}\,k\,\dfrac{q}{d^2}\,\vec{\jmath}$

Noteu que la notació vectorial dependrà del vèrtex triat. No cal que l'estudiant doni el
resultat en notació vectorial, el sentit i direcció també es poden indicar gràficament,
com per exemple en l'esquema de forces adjunt. També s'ha de considerar correcte que el
resultat es doni en funció del $\cos(30°)$, no cal substituir el seu valor.
\end{solapartat}

\begin{solapartat}
\punts{0,25} Definim l'energia de formació o energia potencial electrostàtica d'un
sistema com el treball necessari per constituir-lo.

\punts{0,25} Primer cal l'energia per dur la primera càrrega des de l'infinit fins a la
seva ubicació final. Com que no hi ha més càrregues, el camp elèctric extern i la força a
què està sotmesa la càrrega és zero i, per tant, el treball serà nul.
\[ W_1 = 0 \]

\punts{0,25} Per dur la segona càrrega, caldrà superar la força que aplica la primera
sobre la segona, de manera que el treball que caldrà fer serà igual a la variació de
l'energia potencial:
\[ W_2 = k\,\frac{q_1\cdot q_2}{d} = k\,\frac{q^2}{d} \]

\punts{0,25} Per dur la tercera càrrega, caldrà superar la força que apliquen les dues
primeres càrregues sobre la tercera, de manera que el treball que caldrà fer serà igual a
la variació de l'energia potencial:
\[ W_3 = k\,\frac{q_1\cdot q_3}{d} + k\,\frac{q_2\cdot q_3}{d} = 2k\,\frac{q^2}{d} \]

\punts{0,25} Finalment, si sumen tots els termes tenim: $U = 3k\,\dfrac{q^2}{d}$

També es considerarà vàlid que l'estudiant doni directament l'expressió:
\[ U = k\,\frac{q_1\cdot q_2}{d} + k\,\frac{q_1\cdot q_3}{d} + k\,\frac{q_2\cdot q_3}{d} \]
\end{solapartat}
